This section details the software components within the Client and Cloud tiers, defining their specific responsibilities and interactions.

\subsection{Hybrid Client Components}
\begin{itemize}
    \item \textbf{Native Bridge (Capacitor):} The architectural boundary between the web code and the native OS. It exposes a unified API to access hardware features. It ensures that the \textbf{Geolocation} and \textbf{Accelerometer} continue to operate reliably even when the application is in the background, bridging the gap between the web runtime and the device hardware.
    
    \item \textbf{Sensor Controller \& Filter:} Responsible for logic related to data acquisition.
    \begin{itemize}
        \item It interfaces with the \textbf{Native Bridge} to subscribe to GPS and Accelerometer streams.
        \item It implements a \textbf{Kalman Filter} to smooth raw GPS coordinates.
        \item It analyzes accelerometer data to detect vertical shocks (road surface irregularities like potholes) and buffers this data locally during activity recording.
    \end{itemize}
    
    \item \textbf{Local Persistence Manager:} Implements the "Offline-First" capability. It uses an embedded database (e.g., SQLite via Capacitor plugin) to store activity traces and pending reports when the network is unavailable. It includes a \textbf{Sync Engine} that monitors network reachability and uploads buffered data to the cloud once connectivity is restored.
    
    \item \textbf{BaaS Client SDK:} The abstraction layer provided by the BaaS platform. It manages the persistent WebSocket connection for real-time updates, handles secure API calls to the database, and automatically refreshes authentication tokens.
    
    \item \textbf{Map \& UI Controller:} Manages the visualization of map tiles and path geometries. It handles user inputs (forms, search queries) and renders the visual feedback for the navigation system.
\end{itemize}

\subsection{Managed Cloud Components (BaaS)}
\begin{itemize}
    \item \textbf{Auth Service (Identity Provider):} Manages user lifecycles (sign-up, login, password recovery). It integrates with third-party OAuth providers (e.g., Google, Facebook) and issues secure JWTs containing user roles and IDs.
    
    \item \textbf{Spatial Database (PostgreSQL + PostGIS):} The central data repository. It is an active component that:
    \begin{itemize}
        \item Enforces \textbf{RLS Policies} to secure data access.
        \item Executes complex \textbf{Spatial Queries} (e.g., \texttt{ST\_DWithin}, \texttt{ST\_Intersects}) natively, allowing clients to efficiently retrieve paths within a specific viewport.
        \item Triggers database events (Webhooks) to activate Edge Functions upon data changes.
    \end{itemize}
    
    \item \textbf{Object Storage:} A specialized file storage system for large binary assets, such as user profile pictures or images attached to path reports.
\end{itemize}

\subsection{Edge Functions}
These are stateless functions triggered by database events or scheduled cron jobs:
\begin{itemize}
    \item \texttt{func\_path\_merging}
    \begin{itemize}
        \item \textbf{Trigger:} Database event (INSERT on \texttt{Activities} or \texttt{Feedbacks} table).
        \item \textbf{Responsibility:} It implements the \textbf{Data Consolidation \& Consensus Logic}. When new data (an activity trace or a manual report) is associated with a path, this function aggregates it with existing reports to determine the path's \texttt{CurrentCondition}. It applies a \textbf{Majority Voting Algorithm} weighted by data freshness (e.g., resolving a conflict where 3 recent reports say "Good" and 1 says "Bad" by marking the segment as "Good") to ensure the status reflects the actual road situation.
    \end{itemize}

    \item \texttt{func\_path\_search}
    \begin{itemize}
        \item \textbf{Trigger:} Direct Client Request.
        \item \textbf{Responsibility:} It implements the \textbf{Search \& Ranking Algorithm} defined in the RASD. When a user queries for a route (specifying Origin and Destination), this function:
        \begin{itemize}
            \item Performs a spatial query to find candidate paths connecting the points.
            \item Dynamically calculates a \textbf{Suitability Score} for each candidate. This score combines the \textbf{Geospatial match} (how well it fits the requested Start/End points) with the \textbf{Path Quality} (derived from the consolidated feedback and applying a time-decay factor to older data).
            \item Returns the list of paths ordered by this score.
        \end{itemize}
    \end{itemize}

    \item \texttt{func\_maintenance\_job}
    \begin{itemize}
        \item \textbf{Trigger:} Scheduled Cron Job (e.g., Nightly).
        \item \textbf{Responsibility:} Periodically applies the \textbf{Time-Decay Function} to the database. It archives or lowers the weight of very old reports/feedbacks to ensure that the \texttt{func\_path\_search} always prioritizes fresh information without needing to process the entire history during every query.
    \end{itemize}

    \item \texttt{func\_weather\_enricher}
    \begin{itemize}
        \item \textbf{Trigger:} Database event (INSERT on \texttt{Activities}).
        \item \textbf{Responsibility:} It securely calls external Meteorological APIs (using backend-only API keys) to fetch weather conditions for the specific time and location of the activity, enriching the path metadata.
    \end{itemize}

    \item \texttt{func\_path\_deduplication}
    \begin{itemize}
        \item \textbf{Trigger:} Direct Client Request (RPC via check\_path\_duplicate)
        \item \textbf{Responsibility:} It implements the \textbf{Geometric Integrity Control Logic.} 
        Before a new manual path is saved to the database, this function receives the geometry drawn by the user and executes a high-performance spatial query to verify the existence of already present public paths that are geometrically identical or extremely similar. It returns the validation result to the client, preventing the creation of redundant duplicates and ensuring path uniqueness within the system.
    \end{itemize}
\end{itemize}

% INSERT DIAGRAM 2.2 HERE
% \begin{figure}[H]
% \centering
% \includegraphics[width=\textwidth]{images/diagram_2_2_components.png}
% \caption{Component View}
% \label{fig:component_view}
% \end{figure}